\bcpchapters{Publick Baptism of Infants}{The Ministration of Publick Baptism of Infants}

\small

\begin{leftbar}

\begin{rubric}
The Ministers of every Parish shall often admonish the people that they bring their children to Baptism as soon as possible after birth, and that they defer not the baptism longer than the fourth, or at furthest the fifth, Sunday unless upon a great and reasonable cause.
\end{rubric}

\begin{rubric}
It is desirable where possible that Baptism should be administered upon Sundays and other Holy-days, when the most number of people come together; as well for that the congregation there present may testify the receiving of them that be newly baptized into the number of Christ's Church; as also because in the Baptism of infants every man present may be put in remembrance of his own profession made to God in his Baptism: for which cause it is expedient that Baptism be ministered in the vulgar tongue.
Nevertheless (for sufficient cause), children may be baptized upon any other day.
\end{rubric}

\begin{rubric}
And note, that there shall be for every male-child to be baptized two Godfathers and one Godmother; and for every female, one Godfather and two Godmothers. Nevertheless, when three sponsors cannot conveniently be had, one Godfather and one Godmother shall suffice.
Parents, if need so require, may be sponsors for their own child provided that the child have one other sponsor.
No person shall be admitted to be a sponsor who hath not been baptized.
\end{rubric}

\begin{rubric}
In the absence of the Priest it is lawful that a Deacon baptize infants.
\end{rubric}

\begin{rubric}
When there are children to be baptized, the Parents shall give due notice thereof to the Minister of the Parish.
He shall thereupon appoint the time for the Baptism.
Then the Godfathers and Godmothers, and the people with the children, must be ready at the Font either immediately after the last Lesson, or after the Third Collect at Morning or Evening Prayer; or at such other time as he in his discretion shall think fir.
\end{rubric}

\end{leftbar}

\begin{rubric}
It is certain by God's Word, that children which are baptized, dying before they commit actual sin, are undoubtedly saved.
\end{rubric}

\begin{rubric}
To take away all scruple concerning the use of the sign of the Cross in Baptism; the true explication thereof, and the just reasons for the retaining of it, may be seen in the xxxth Canon, first published in the Year MDCIV.
\end{rubric}

\normalsize

\newpage

\begin{rubric}
And the Priest coming to the Font, (which is then to be filled with pure water,) and standing there shall say,
\end{rubric}

\lett{H}{ath} this Child been already baptized, or no?
\\

\begin{rubric}
If they answer, No: then shall the Priest proceed as followeth:
\end{rubric}

\lett{D}{early} beloved, forasmuch as all men are conceived and born in sin: and that our Saviour Christ saith, None can enter into the kingdom of God, except he be regenerate and born anew of Water and of the Holy Ghost: I beseech you to call upon God the Father, through our Lord Jesus Christ, that of his bounteous mercy he will grant to this Child that thing which by nature he cannot have; that he may be baptized with Water and the Holy Ghost, and received into Christ's holy Church, and be made a lively member of the same.

\begin{rubric}
Then shall the Priest say,
\end{rubric}

\lup{Let us pray.}

\lett{A}{lmighty} and everlasting God, who of thy great mercy didst save Noah and his family in the ark from perishing by water; and also didst safely lead the children of Israel thy people through the Red Sea, figuring thereby thy holy Baptism; and by the Baptism of thy well-beloved Son Jesus Christ, in the river Jordan, didst sanctify Water to the mystical washing away of sin: We beseech thee, for thine infinite mercies, that thou wilt mercifully look upon this Child; wash him and sanctify him with the Holy Ghost; that he, being delivered from thy wrath, may be received into the ark of Christ's Church; and being stedfast in faith, joyful through hope, and rooted in charity, may so pass the waves of this troublesome world, that finally he may come to the land of everlasting life, there to reign with thee world without end; through Jesus Christ our Lord.
\inlineresp{Amen.}

\lett{A}{lmighty} and immortal God, the aid of all that need, the helper of all that flee to thee for succour, the life of them that believe, and the resurrection of the dead: We call upon thee for this Infant, that he, coming to thy holy Baptism, may receive remission of his sins by spiritual regeneration. Receive him, 0 Lord, as thou hast promised by thy well-beloved Son, saying, Ask, and ye shall have; seek, and ye shall find; knock, and it shall be opened unto you: So give now unto us that ask; let us that seek find; open the gate unto us that knock; that this Infant may enjoy the everlasting benediction of thy heavenly washing, and may come to the eternal kingdom which thou hast promised by Christ our Lord.
\inlineresp{Amen.}

\begin{rubric}
Then shall the people stand up, and the Priest shall say,
\end{rubric}

\begin{center}
Hear the words of the Gospel, written by Saint Mark,\\
in the tenth chapter, at the thirteenth verse
\end{center}

\lett{T}{hey} brought young children to him, that he should touch them: and his disciples rebuked those that brought them.
But when Jesus saw it, he was much displeased, and said unto them, Suffer the little children to come unto me, and forbid them not: for of such is the kingdom of God.
Verily I say unto you, Whosoever shall not receive the kingdom of God as a little child, he shall not enter therein.
And he took them up in his arms, put his hands upon them, and blessed them.

\begin{rubric}
After the Gospel is read, the Minister shall make this brief Exhortation upon the words of the Gospel.
\end{rubric}

\lett{B}{eloved}, ye hear in this Gospel the words of our Saviour Christ, that he commanded the children to be brought unto him; how he blamed those that would have kept them from him; how he exhorteth all men to follow their innocency. Ye perceive how by his outward gesture and deed he declared his good will toward them; for he embraced them in his arms, he laid his hands upon them, and blessed them. Doubt ye not therefore, but earnestly believe, that he will likewise favourably receive this present Infant; that he will embrace him with the arms of his mercy; that he will give unto him the blessing of eternal life, and make him partaker of his everlasting kingdom. Wherefore we being thus persuaded of the good will of our heavenly Father towards this Infant, declared by his Son Jesus Christ; and nothing doubting but that he favourably alloweth this charitable work of ours in bringing this Infant to his holy Baptism; let us faithfully and devoutly give thanks unto him, and say,

\lett{A}{lmighty} and everlasting God, heavenly Father, we give thee humble thanks, for that thou hast vouchsafed to call us to the knowledge of thy grace, and faith in thee: Increase this knowledge, and confirm this faith in us evermore. Give thy Holy Spirit to this Infant, that he may be born again, and be made an heir of everlasting salvation; through our Lord Jesus Christ, who liveth and reigneth with thee and the Holy Spirit, now and for ever.
\inlineresp{Amen.}

\begin{rubric}
Then shall the Priest speak unto the Godfathers and Godmothers on this wise.
\end{rubric}

\lett{D}{early} beloved, ye have brought this Child here to be baptized, ye have prayed that our Lord Jesus Christ would vouchsafe to receive him, to release him of his sins, to sanctify him with the Holy Ghost, to give him the kingdom of heaven, and everlasting life. Ye have heard also that our Lord Jesus Christ hath promised in his Gospel to grant all these things that ye have prayed for: which promise he, for his part, will most surely keep and perform. Wherefore, after this promise made by Christ, this Infant must also faithfully, for his part, promise by you that are his sureties, (until he come of age to take it upon himself,) that he will renounce the devil and all his works, and constantly believe God's holy Word, and obediently keep his commandments.
\\

\lup{I demand therefore,}

\lett{D}{ost} thou, in the name of this Child, renounce the devil and all his works, the vain pomp and glory of the world, with all covetous desires of the same, and the carnal desires of the flesh, so that thou wilt not follow, nor be led by them?

\inrub{Answer.} I renounce them all.

\lett{D}{ost} thou believe in God the Father Almighty, Maker of heaven and earth?
And in Jesus Christ his only-begotten Son our Lord? And that he was conceived by the Holy Ghost; born of the Virgin Mary; that he suffered under Pontius Pilate, was crucified, dead, and buried; that he went down into hell, and also did rise again the third day; that he ascended into heaven, and sitteth at the right hand of God the Father Almighty; and from thence shall come again at the end of the world, to judge the quick and the dead?
And dost thou believe in the Holy Ghost; the holy Catholick Church; the Communion of Saints; the Remission of sins; the Resurrection of the flesh; and everlasting life after death?

\inrub{Answer.} All this I stedfastly believe.

\lett{W}{ilt} thou be baptized in this faith?

\inrub{Answer.} That is my desire.

\lett{W}{ilt} thou then obediently keep God's holy will and commandments, and walk in the same all the days of thy life?

\inrub{Answer.} I will.

\lett{O}{ merciful} God, grant that the old Adam in this Child may be so buried, that the new man may be raised up in him. \inlineresp{Amen.}\\
Grant that all carnal affections may die in him, and that all things belonging to the Spirit may live and grow in him. \inlineresp{Amen.}\\
Grant that he may have power and strength to have victory, and to triumph, against the devil, the world, and the flesh. \inlineresp{Amen.}\\
Grant that whosoever is here dedicated to thee by our office and ministry may also be endued with heavenly virtues, and everlastingly rewarded, through thy mercy, O blessed Lord God, who dost live, and govern all things, world without end. \inlineresp{Amen.}

\lett{A}{lmighty}, everliving God, whose most dearly beloved Son Jesus Christ, for the forgiveness of our sins, did shed out of his most precious side both water and blood; and gave commandment to his disciples, that they should go teach all nations, and baptize them In the Name of the Father, and of the Son, and of the Holy Ghost: Regard, we beseech thee, the supplications of thy congregation; sanctify this Water to the mystical washing away of sin; and grant that this Child, now to be baptized therein, may receive the fulness of thy grace, and ever remain in the number of thy faithful and elect children; through Jesus Christ our Lord. \inlineresp{Amen.}

